%\usepackage[utf8]{inputenc}
\usepackage[british]{babel}
\usepackage{relsize,setspace}
%\usepackage{stmaryrd}
\usepackage{url} % [hyphens] passed in short.tex/long.tex
%\usepackage[font=it]{caption}
\usepackage{tikz,pgf}\usetikzlibrary{shapes,backgrounds}
\usepackage{amsmath,amssymb}
\usepackage[amsmath]{ntheorem}
%\usepackage[pdftex,colorlinks=true,linkcolor=black,citecolor=black,menucolor=black,urlcolor=black]{hyperref}
\usepackage{cleveref}
\usepackage{multirow}
\usepackage{comment}
\usepackage{xspace}
%\usepackage{pdflscape}
%\usepackage{yhmath}
\usepackage{longtable}
\usepackage{stmaryrd}

% environments
%\theoremstyle{plain}
\theoremseparator{.}
\theoremheaderfont{\normalfont\bfseries}
\theorembodyfont{\slshape}
\newtheorem{theorem}{Theorem}
\newtheorem{lemma}[theorem]{Lemma}
\newtheorem{proposition}[theorem]{Proposition}
\newtheorem{corollary}[theorem]{Corollary}
\newtheorem{observation}[theorem]{Observation}

\theorembodyfont{\upshape}
\theoremsymbol{\ensuremath{\Diamond}}
\newtheorem{definition}{Definition}
\newtheorem{example}{Example}
\theoremstyle{nonumberplain}
\theoremheaderfont{\normalfont\bfseries}
\theorembodyfont{\slshape}
\newtheorem{claim}{Claim}
\newtheorem{nonumberlemma}{Lemma}
\newtheorem{nonumbertheorem}{Theorem}
\theoremstyle{nonumberplain}
\theoremheaderfont{\normalfont\itshape}
\theoremsymbol{\ensuremath{\Box}}
\newtheorem{proof}{Proof}
\newtheorem{proofsketch}{Proof sketch} % note that this does not care \iflong
% this is for nesting proofs, e.g. proving some distinguished sub-claim within a proof
\theoremstyle{nonumberplain}
\theoremheaderfont{\normalfont\itshape}
\theoremsymbol{\ensuremath{\Diamond}}
\newtheorem{claimproof}{Proof of the claim}
\newtheorem{remark}{Remark}

%% for producing a long version of the paper with proofs
\usepackage{ifthen}
\usepackage{environ}
%% this could in principle be used to use “long” and “short” as non-exclusive or non-exhaustive, or both
%\newif\ifshort
\newif\iflong
% use shortproofs for proofs that should appear in the short paper, but not in the long paper
\NewEnviron{shortproof}{\iflong\else\expandafter\begin{proof}\BODY\end{proof}\fi}
% shortproofsketches are a variant (special case) of short proofs where the environment says “Proof sketch” instead of “Proof”
\NewEnviron{shortproofsketch}{\iflong\else\expandafter\begin{proofsketch}\BODY\end{proofsketch}\fi}
% use longproofs for proofs that should appear in the long paper only (i.e. not in the short paper)
\NewEnviron{longproof}{\iflong\expandafter\begin{proof}\BODY\end{proof}\fi}

\usepackage{todonotes}
\usepackage{xcolor}
\definecolor{lightblue}{HTML}{55AAFF}
\definecolor{lightpurple}{HTML}{DD77BB}
\definecolor{lightgray}{HTML}{DDDDDD}
\definecolor{strongyellow}{HTML}{FFE219}
\definecolor{darkorange}{HTML}{C2570C}

\definecolor{codegreen}{rgb}{0,0.6,0}
\definecolor{codegray}{rgb}{0.5,0.5,0.5}
\definecolor{codepurple}{rgb}{0.58,0,0.82}
\definecolor{backcolour}{rgb}{0.95,0.95,0.92}

\newcommand{\hstodo}[2][]{\todo[author=HS,color=teal!50,#1]{#2}}
\newcommand{\drtodo}[2][]{\todo[author=DR,color=strongyellow,#1]{#2}}
\newcommand{\sgtodo}[2][]{\todo[author=SG,color=lightblue,#1]{#2}}
\newcommand{\pktodo}[2][]{\todo[author=PK,color=lightpurple,#1]{#2}}
\newcommand{\jhtodo}[2][]{\todo[author=PK,color=lightgray,#1]{#2}}

\newcommand{\displaytodo}[1]{#1}
\newcommand{\rem}[1]{\displaytodo{\todo[inline,color=violet!40!white]{\textbf{NOTE:} #1}}}
\newcommand{\tod}[1]{\displaytodo{\todo[inline,color=red!40]{\textbf{TODO:} #1}}}

% LaTeX macros
\newcommand{\ghost}[1]{\raisebox{0pt}[0pt][0pt]{\makebox[0pt][l]{#1}}}
% math & logic
\newcommand{\N}{\mathbb{N}}
\newcommand{\Z}{\mathbb{Z}}
\newcommand{\R}{\mathbb{R}}
% number theory
%\newcommand{\divides}{\mid}
% why is this not a standard command but \gcd is?
%\newcommand{\lcm}{\mathbin{\mathrm{lcm}}}
% renew gcd to make it look like lcm
%\renewcommand{\gcd}{\mathbin{\mathrm{gcd}}}
% set theory
\newcommand{\set}[1]{\left\{#1\right\}}
\newcommand{\pset}[1]{2^{#1}}%
\newcommand{\sset}[1]{\left\{\,#1\,\right\}}
\newcommand{\guard}{\ \middle\vert\ }
\newcommand{\card}[1]{\left\vert#1\right\vert}
%\newcommand{\size}[1]{\left\Vert#1\right\Vert}
\newcommand{\tuple}[1]{\left(#1\right)}
%\newcommand{\tand}{\mathrel{\:\:\&\:\:}}
\newcommand{\tand}{\;\text{and}\;}
%\newcommand{\timplies}{\implies}
\newcommand{\timplies}{\quad\text{implies}\quad}
\newcommand{\define}[1]{\emph{#1}}
% use next command as define, but if you want something else (#2) to occur in the index/glossary
\newcommand{\defineas}[2]{\emph{#1}}
% thin spaces
\newcommand{\ts}{\mkern1mu}
\newcommand{\nts}{\mkern-1mu}
% ordinal numbers
\newcommand{\ordn}{\alpha}
\newcommand{\ordnalt}{\alpha'}
\newcommand{\ordsucc}{\mathord{\ordn+1}}
\newcommand{\ordlim}{\beta}
%% meta-level logic
% notation for meta-level quantification
%\newcommand{\qexists}[2]{(\exists#1)#2}
%\newcommand{\qforall}[2]{(\forall#1)#2}
%\newcommand{\onlyif}{\Longleftarrow}
% object-level macros
\newcommand{\eqdef}{\mathrel{:=}}
\newcommand{\iffdef}{\mathrel{\;:\!\!\iff\,}}
%% ebnf
% production rule/definition
\newcommand{\ebnfeq}{\mathrel{::=}}
% alternative
\newcommand{\ebnfalt}{\mathbin{\mid}}



% Define typographic struts, as suggested by Claudio Beccari
%   in an article in TeX and TUG News, Vol. 2, 1993.
\newcommand\Tstrut{\rule{0pt}{2.6ex}}       % = `top' strut
\newcommand\Bstrut{\rule[-1ex]{0pt}{0pt}}   % = `bottom' strut
\newcommand\Strut{\rule[-1ex]{0pt}{2.6ex}}  % = `both' strut

% Down right arrow
%\newcommand{\drarrow}{\rotatebox[origin=c]{180}{$\Lsh$}}







%%%%%%


%\newcommand{\citet}[1]{\citeauthor{#1}~[\citeyear{#1}]}

\newcommand{\explain}[1]{\ensuremath{#1}\dotfill\texttt{\detokenize{#1}}}

% explain macro and usage as table line
\newcommand{\taex}[2]{\ensuremath{#1}&\texttt{\detokenize{#1}}&#2\\}
\newcommand{\taexn}[2]{\textcolor{blue}{\ensuremath{#1}}&\textcolor{blue}{\texttt{\detokenize{#1}}}&\textcolor{blue}{#2}\\}
\newcommand{\taexp}[2]{\textcolor{purple}{\ensuremath{#1}}&\textcolor{purple}{\texttt{\detokenize{#1}}}&\textcolor{purple}{#2}\\}
\newcommand{\taexheading}[1]{\hline\multicolumn{3}{c}{\textbf{#1}}\\\hline}

\newlength{\narrowgatherskip}
\newenvironment{narrowgather}[1][2pt]{%
    \setlength{\narrowgatherskip}{#1}%
    \vskip\narrowgatherskip\mbox{}\hfill\(}{\)\hfill\mbox{}\vskip\narrowgatherskip\noindent}


%% Formal Language
\newcommand{\propL}{\mathfrak{L}}
%% Fuzzy Universe 
%\newcommand{\fuzzyUniverse}[1]{#1}
%% Fuzzy Connectives 
\newcommand{\flpcon}{f}
%% Family of Fuzzy Connectives
\newcommand{\Flpcon}{\mathcal{F}}
%% atoms in general (propositional vocabulary)
\newcommand{\Atoms}{\Pi}
\newcommand{\Atomsalt}{\Pi'}
\newcommand{\atom}{p}
\newcommand{\litset}{\alpha}
\newcommand{\litsetsup}[2][i]{\litset^{(#1,#2)}}
\newcommand{\litsetalt}{\beta}
\newcommand{\litsetaltsup}[2][i]{\litsetalt^{(#1,#2)}}
%% Atoms occurring in program
%\newcommand{\flpatom}[1][\flp]{\Pi_{#1}}%
\newcommand{\flpatom}{\Atoms}%
%% Fuzzy Interpretations 
\newcommand{\fint}{I}
\newcommand{\fintsup}[2][i]{\fint^{(#1,#2)}}
\newcommand{\fintalt}{J}
\newcommand{\fintaltsup}[2][i]{\fintalt^{(#1,#2)}}
\newcommand{\fintset}{\mathcal{J}}
\newcommand{\lint}{L}
\newcommand{\lintsup}[1][i]{\lint^{(#1)}}
\newcommand{\uint}{U}
\newcommand{\uintsup}[1][i]{\uint^{(#1)}}
\newcommand{\lintalt}{M}
\newcommand{\uintalt}{T}
\newcommand{\fintpair}{\mathbf{i}}
\newcommand{\fintpairsup}[1][k]{\fintpair^{(#1)}}
\newcommand{\fintpairalt}{\mathbf{j}}
\newcommand{\fintpairaltsup}[1][k]{\fintpairalt^{(#1)}}
%% set of all non-paraconsistent (standard) fuzzy interpretations
\newcommand{\nint}{K}
\newcommand{\nintalt}{L}
\newcommand{\Nint}{\mathcal{\nint}}
%% set of all non-paraconsistent (extended) fuzzy Interpretations
\newcommand{\eint}{G}
\newcommand{\eintalt}{H}
\newcommand{\Eint}{\mathcal{\eint}}
% \newcommand{\Eint}{\mathcal{\eint}^{\flpneg}}
\newcommand{\dom}[1][\eint]{\mathrm{dom}(#1)}
%% Set of all paraconsistent fuzzy Interpretations
\newcommand{\Fint}{\mathcal{\fint}}
%% Corresponding truth value function
\newcommand{\fformeval}[2]{\widehat{#1}\!\left(#2\right)}
\newcommand{\fformnoteval}[2]{#2^{#1}}
\newcommand{\fnnformeval}[2]{\widehat{\widehat{#1}}\!\left(#2\right)}
\newcommand{\fnformeval}[3]{\fnnformeval{(#1,#2)}{#3}}
\newcommand{\flpcontf}[1]{#1^{\text{\raisebox{0.3ex}{\large$\mkern-1mu\mathord{\cdot}$}}}}
\newcommand{\flpcontff}[1]{#1^{\text{\raisebox{0.3ex}{\large$\mkern-1mu\mathord{\cdot\mkern-2mu\cdot}$}}}}
% evaluation for non-paraconsistent interpretations
\newcommand{\nfformeval}[2]{\left[#1\right]\!\left(#2\right)}
\newcommand{\nfnnformeval}[2]{\left[\!\left[#1\right]\!\right]\!\left(#2\right)}
\newcommand{\nfnformeval}[3]{\nfnnformeval{(#1,#2)}{#3}}
% some shorthands
\newcommand{\flpnotf}{\mathord{\flpcontf{\flpnot}\mkern-2mu}}
\newcommand{\flpnotff}{\mathord{\flpcontff{\flpnot}\mkern-2mu}}
\newcommand{\flpnegf}{\mathord{\flpcontf{\flpneg}\mkern-2mu}}
\newcommand{\flpnegff}{\mathord{\flpcontff{\flpneg}\mkern-2mu}}
%% LaTeX macros for fuzzy logic programming
\newcommand{\flpif}[1]{\overset{#1}{\gets}}
\newcommand{\flpnot}{\mathord{\sim}}
\newcommand{\flpneg}{\neg}
\newcommand{\flpand}{\land}
\newcommand{\flpor}{\lor}
\newcommand{\flpagg}{\mathbin{@}}
\newcommand{\flpweight}{\vartheta}
% flp operators 
\newcommand{\flp}{P}
% \newcommand{\flp}{\lp}
\newcommand{\flpoperator}[2]{#1_{#2}}
\newcommand{\flpop}[1][\flp]{\flpoperator{T}{#1}}
\newcommand{\flpap}[1]{\flpoperator{\mathcal{A}}{#1}}
\newcommand{\flpapp}{\flpap{\flp}}
\newcommand{\flpappa}{\fst{\flpapp}}
\newcommand{\flpappb}{\snd{\flpapp}}
%% flp general
\newcommand{\flpr}{\atom \flpif{\flpweight}_i \body}
\newcommand{\flprUW}{\atom \flpif{}_i \body}
\newcommand{\nflpr}{\flpneg \atom \flpif{\flpweight}_i \body}
\newcommand{\nflprUW}{\flpneg \atom \flpif{}_i \body}
\newcommand{\lflpr}{\literal \flpif{}_i \body}
\newcommand{\flprRed}[1][\fint]{\atom \flpif{\flpweight}_i \body_{#1}}
\newcommand{\eflpr}{\literal\lpif{}\body}
\newcommand{\eflprp}{\atom\lpif{}\body}
\newcommand{\eflprn}{\flpneg\atom\lpif{}\body}
\newcommand{\body}{\mathcal{B}}
\newcommand{\pbody}{\body^{\not\flpnot}}
\newcommand{\nbody}{\body^{\flpnot}}
\newcommand{\head}{\mathcal{H}}
\newcommand{\bodies}[1][\literal]{\mathbf{B}_{#1}}


\newcommand{\flpcontfStraccia}[1]{#1^{\text{\raisebox{0.3ex}{\large$\mathord{\cdot}\mathord{\cdot}$}}}}
\newcommand{\TpStraccia}[1][\flp]{\mathcal{S}_{#1}} % S for Straccia
\newcommand{\FpStraccia}[2][\flp]{F_{#1,#2}}
\newcommand{\SpStraccia}{s}
\newcommand{\WFStraccia}{\mathcal{AW}}

% Closed World Assumption
\newcommand{\cwa}{CWA}

%Abbreviation for various "fuzzy logic program(s)" when needed
\newcommand{\flptext}{FLP\xspace}
\newcommand{\flpstext}{FLPs\xspace}
\newcommand{\pflptext}{p\flptext}
\newcommand{\pflpstext}{p\flpstext}
\newcommand{\nflptext}{n\flptext}
\newcommand{\nflpstext}{n\flpstext}
\newcommand{\eflptext}{e\flptext}
\newcommand{\eflpstext}{e\flpstext}

% normal non-paraconsistent logic program
\newcommand{\nflp}{Q}

% modus ponens inference candidate
\newcommand{\mopo}{\mathit{mp}}

% example program atoms
\newcommand{\epred}[1]{\mathtt{#1}}
\newcommand{\drive}{\epred{drive}}
\newcommand{\obstacle}{\epred{obstacle}}
\newcommand{\camera}{\epred{camera}}
\newcommand{\proximity}{\epred{proximity}}

% Saad syntax
\newcommand{\ann}[2]{#1\mathord{:}#2}
\newcommand{\gecon}[1]{\mathord{\trianglerighteq_{#1}}}
\newcommand{\gtcon}[1][0]{\mathord{\triangleright_{#1}}}

% Saad fixpoint operator
\newcommand{\Saadop}{O}
\newcommand{\SaadP}{Q}

% Reduct of Saad rule
\newcommand{\Saadredr}[1][\eint]{r_{#1}}
% Reduct of Saad program
\newcommand{\Saadredp}[1][\eint]{\SaadP_{#1}}

% Map from literal interpretations to pairs
\newcommand{\litTolit}[1][\eint]{\fint_{#1}}

% remainder of the head after removing p
\newcommand{\headminus}[1][\atom]{\head_{#1}}
%% LaTeX macros for logic programming
\newcommand{\lpr}{r}
\newcommand{\hr}[1]{H(#1)}
\newcommand{\br}[1]{B(#1)}
\newcommand{\pbr}[1]{B(#1)^+}
\newcommand{\nbr}[1]{B(#1)^-}
\newcommand{\lpif}{\gets}
\newcommand{\lpand}{\land}
%\newcommand{\lpnot}{\neg}
\newcommand{\lpnot}{\mathord{\sim}}
\newcommand{\lpatoms}[1]{\mathcal{A}(#1)}
\newcommand{\lits}[1]{\mathcal{L}(#1)}
\newcommand{\compl}[1]{\mathcal{C}(#1)}
\newcommand{\aset}{\mathfrak{S}}


% literals
\newcommand{\literals}[1]{#1^{\pm}}
\newcommand{\dnliterals}[1]{#1^{\pm-}}
\newcommand{\nliterals}[1]{#1^{-}}
\newcommand{\literal}{\ell}
% example logic program
\newcommand{\elp}{P}



\newcommand{\A}[1][\lp]{\mathcal{A}(#1)}
\newcommand{\AS}[1][\lp]{\mathcal{AS}(#1)}                    %
\newcommand{\BC}[1][\lp]{\mathcal{BC}(#1)}                    %
\newcommand{\CC}[1][\lp]{\mathcal{CC}(#1)}                    %
\newcommand{\F}[1][\lp]{\mathcal{F}(#1)}                      %
\newcommand{\FI}[1][\lp]{\mathcal{F}^+(#1)}                   %
\newcommand{\FE}[1][\lp]{\mathcal{F}^-(#1)}                   %
\newcommand{\FA}[1][\lp]{\mathcal{F}_{a}(#1)}                 %
\newcommand{\SA}[2][\lp]{S^{\models_{#1} #2}}                 %
\newcommand{\NSA}[2][\lp]{S^{\not \models_{#1} #2}}           %
\newcommand{\C}[2]{\mathcal{C}(#1, #2)}                       %
\newcommand{\MC}[2]{\mathcal{MC}(#1, #2)}                     %
\newcommand{\DF}[1][\lp]{\mathcal{DF}(#1)}                    %

\newcommand{\trans}[3][\lp]{#1 \overset{#2}{\Rightarrow} #3}  %

\newcommand{\route}[1][]{\langle #1 \rangle}                  %
\newcommand{\De}[2][\lp]{\Delta^{#1}_{#2}}                    %
\newcommand{\RE}[2][f]{\mathcal{R}(#2, #1)}                   %
\newcommand{\MRE}[2][f]{\mathcal{MR}(#2, #1)}                 %

\newcommand{\vgo}[1][]{\nu_{\mathit{go}}#1}                   %
\newcommand{\vfr}[1][]{\nu_{\mathit{free}}#1}                 %
\newcommand{\pace}[2]{\mathcal{P}_{#1}(#2)}             %

%% LaTeX macros for lattice theory
% lattice
\newcommand{\lset}{V}
\newcommand{\lelem}{\alpha}
\newcommand{\lelemalt}{\beta}
%\newcommand{\latset}{\lset}
\newcommand{\lleq}{\leqslant}
\newcommand{\lgeq}{\geqslant}
\newcommand{\llt}{<}
\newcommand{\lgt}{>}
\newcommand{\lub}{\vee}
\newcommand{\glb}{\wedge}
\newcommand{\biglub}{\bigvee\!}
\newcommand{\bigglb}{\bigwedge\!}
\newcommand{\least}{0}
\newcommand{\grtst}{1}
\newcommand{\upperb}[1]{#1^u}
\newcommand{\lowerb}[1]{#1^\ell}
\newcommand{\upset}[1]{\uparrow#1}
\newcommand{\downset}[1]{\downarrow#1}
\newcommand{\parset}{X}

% least and greatest elements in bilattice
\newcommand{\ibot}{\perp}
\newcommand{\tbot}{\texttt{f}}
\newcommand{\itop}{\top}
\newcommand{\ttop}{\texttt{t}}

%% LaTeX macros for approximation fixpoint theory
% components of pairs: prefix command
\newcommand{\fst}[1]{#1_1}
\newcommand{\snd}[1]{#1_2}
% postfix command
\newcommand{\first}{'}
\newcommand{\second}{''}
% operators
\newcommand{\op}{O}
\newcommand{\ap}{\mathcal{A}}
\newcommand{\opap}{\mathcal{O}}
\newcommand{\apa}{\fst{\ap}}
\newcommand{\apb}{\snd{\ap}}
\newcommand{\stable}[1]{#1^{\mathit{st}}}
\newcommand{\stap}{\stable{\ap}}
\newcommand{\opst}{\stable{\opap}}
\newcommand{\complete}[1]{\mathord{c\stable{#1}}}
% lp operators
\newcommand{\lpoperator}[2]{#1_{#2}}
\newcommand{\lp}{P}
\newcommand{\lpop}[1][\lp]{\lpoperator{T}{#1}}
\newcommand{\lpopp}{\lpop{\lp}}
\newcommand{\lpap}[1][\lp]{\lpoperator{\mathcal{T}}{#1}}
\newcommand{\lpapp}{\lpap{\lp}}
\newcommand{\lpappa}{\fst{\lpapp}}
\newcommand{\lpappb}{\snd{\lpapp}}
% af operators
\newcommand{\afop}[1]{\lpoperator{U}{#1}}
\newcommand{\af}{F}
\newcommand{\afopp}{\afop{\af}}
\newcommand{\afap}[1]{\lpoperator{\mathcal{U}}{#1}}
\newcommand{\afapp}{\afap{\af}}
% bilattice for a lattice carrier set
\newcommand{\bil}[1][\lset]{#1^2}
% consistent bilattice
\newcommand{\bic}[1][\lset]{#1^c}
% orderings
\newcommand{\ileq}{\leq_i}
\newcommand{\ilt}{<_i}
\newcommand{\pleq}{\leq_p}
\newcommand{\plt}{<_p}
\newcommand{\plub}{\oplus}
\newcommand{\pglb}{\otimes}
\newcommand{\bigplub}{\bigoplus}
\newcommand{\tleq}{\leq_t}
\newcommand{\tlt}{<_t}
\newcommand{\tlub}{\lor_t}
\newcommand{\tglb}{\land_t}
\newcommand{\bigtlub}{\bigvee_t}
\newcommand{\bigtglb}{\bigwedge_t}
% fixpoints
\newcommand{\lfp}{\mathit{lfp}}
\newcommand{\gfp}{\mathit{gfp}}
% another operator for comparison
\newcommand{\bp}{B}
\newcommand{\ult}[1]{\mathcal{U}_#1}
\newcommand{\ultap}{\mathcal{U}}
\newcommand{\ultimate}[1]{\mathord{\ultap_#1}}
\newcommand{\ultop}{\ult{\op}}
\newcommand{\ultlpop}{\ult{\flp}}
% operator precision ordering
\newcommand{\opleq}{\leq_{\mathit{oi}}}

% specific operator for logic programs (supported models)
\newcommand{\lpsu}[1]{\lpop{#1}}
\newcommand{\lpst}[1]{\stable{\lpap{#1}}}



%% propositional logic
\newcommand{\Atom}{\mathcal{A}}%
\newcommand{\formula}{\varphi}%
\newcommand{\formulaalt}{\psi}%
\newcommand{\sub}[1]{\mathrm{sub}'(#1)}%
\newcommand{\subi}[1]{\mathrm{sub}(#1)}%
\newcommand{\bigland}{\bigwedge}
\newcommand{\biglor}{\bigvee}
\newcommand{\limplies}{\rightarrow}
\newcommand{\lonlyif}{\leftarrow}
\newcommand{\liff}{\leftrightarrow}
\newcommand{\lxor}{\oplus} %{\not\leftrightarrow}
% propositional language
\newcommand{\lpl}{\mathfrak{L}}%
% first-order logic
\newcommand{\svec}[1]{\vec{#1}\,}
% prototypical quantifier
\newcommand{\Q}{\rotatebox[origin=c]{180}{\scalebox{0.90}{\raisebox{-1.6ex}[0ex][0ex]{$\mathsf{Q}$}}}}
%% distinguishing model relation and entailment relation, as they have different types: (Interpretation, Formula) vs. (Set of Formula, Formula)
\newcommand{\entails}{\models}
% truth values
\def\tvt{\mathbf{t}}
\def\tvf{\mathbf{f}}
\def\two{\set{\tvt,\tvf}}
\def\tvu{\mathbf{u}}
\def\three{\set{\tvt,\tvf,\tvu}}
\def\tvi{\mathbf{i}}
\def\four{\set{\tvt,\tvf,\tvu,\tvi}}
%% unit interval 
\newcommand{\I}{[0;1]}
% non-monotonic reasoning
\newcommand{\dentails}{\mathbin{\mid\mathrel{\mkern-3.6mu}\approx}}
%% fuzzy logic
\newcommand{\flnot}{\mathord{\sim}}
\newcommand{\flneg}{\neg}
\newcommand{\fland}{\land}
\newcommand{\flor}{\lor}
\newcommand{\bigfland}{\bigwedge}
\newcommand{\bigflor}{\bigvee}
\newcommand{\flif}{\rightarrow}
\newcommand{\flonlyif}{\leftarrow}
\newcommand{\fliff}{\leftrightarrow}
\newcommand{\flcon}{f}
\newcommand{\Flcon}{\mathcal{F}}
\newcommand{\flcontf}[1]{\mathbin{\dot{#1}}}
\newcommand{\flcontfp}[1]{\mathbin{\ddot{#1}}}
\newcommand{\finttf}[1]{\hat{#1}}
\newcommand{\flnottf}{\mathord{\flcontf{\flnot}}}
\newcommand{\flnegtf}{\mathord{\flcontf{\flneg}}}
\newcommand{\iset}{K}






\newcommand{\lnn}{\mathrm{lnn}}
\DeclareMathOperator{\upo}{\ensuremath{\uparrow}}%
\DeclareMathOperator{\Upo}{\ensuremath{\Uparrow}}%
\DeclareMathOperator{\dpo}{\ensuremath{\downarrow}}%
\DeclareMathOperator{\Dpo}{\ensuremath{\Downarrow}}%

\newcommand{\lowb}[1]{\underline{#1}}
\newcommand{\upb}[1]{\overline{#1}}

\newcommand{\Upop}{U}
\newcommand{\Dpop}{D}

% indicate that something is a pair
\newcommand{\pair}[1]{\mathbf{#1}}
% swap elements of a tuple
\newcommand{\rev}[1]{#1^{\triangleleft}}


%%%% Dominiks Macros
% dr
\newcommand{\mino}{\mathrm{min}}%
\newcommand{\maxo}{\mathrm{max}}%
%\DeclareMathOperator{\lneg}{\ensuremath{!}}%
\DeclareMathOperator{\lneg}{\ensuremath{\sim}}%
\newcommand{\signature}{\Sigma}%
%\newcommand{\upo}{\mathcal{U}}%
\newcommand{\aggb}{\rhd}%
\newcommand{\Aggb}{\blacktriangleright}%
% Interpretation on LNNs
\newcommand{\itp}{I}%
% Set of Interpretations on LNNs
\newcommand{\Itp}{\mathcal{I}}%

\newcommand{\inti}[1]{[\![#1]\!]}% 


\newcommand{\dpneg}{\hat{\neg}}
\newcommand{\dpand}{\hat{\land}}
\newcommand{\dpor}{\hat{\lor}}
\newcommand{\dpif}{\hat{\rightarrow}}

\newcommand{\bch}[1]{\left\llbracket #1 \right\rrbracket}

%% Macros for Possibilistic Disjunctive Programs (Nieves2007possibilisitc)

\newcommand{\patom}{a}
\newcommand{\pAtoms}{S}
\newcommand{\pform}{\varphi}
\newcommand{\proj}[1][\patom]{#1^{*}}
\newcommand{\prule}{r}
\newcommand{\pRules}{N}
\newcommand{\pflp}{P}
\newcommand{\Pint}{\mathcal{PS}}
\newcommand{\possleq}{\preceq}
\newcommand{\pnec}{\alpha}

\newcommand{\pint}{\mathcal{A}}
\newcommand{\pintalt}{\mathcal{A}'}

\newcommand{\projrule}[1][\prule]{n(#1)}

\newcommand{\PintToFint}{\zeta}
\newcommand{\possglb}{\sqcap}
\newcommand{\posslub}{\sqcup}

% Reduction operator
\newcommand{\pred}[2]{{#1}^{#2}}

% Inference
\newcommand{\PLinf}{\vdash_{\mathcal{PL}}}

\newcommand{\GGPPE}{\mathcal{G}}
\newcommand{\plpop}{T}
\newcommand{\Semmin}[1]{\mathrm{Sem_{min}}(#1)}
\newcommand{\pFacts}[2]{\mathrm{Facts}(#1,\, #2)}

\newcommand{\pflpalt}{Q}

\newcommand{\PosToFuz}[1]{#1_{trans}}

\newcommand{\dflptext}{dFLP\xspace}
\newcommand{\dflpstext}{dFLPs\xspace}

\newcommand{\gdis}{D}
\newcommand{\gdisalt}{D'}
\newcommand{\gsub}{\sqsubseteq}
\newcommand{\gpair}[2]{\tuple{#1,#2}}
\newcommand{\Gdis}{\mathcal{D}}
\newcommand{\dflpr}{\gdis \flpif{\flpweight} \body}

\newcommand{\gop}[1][\flp]{\flpoperator{\mathcal{C}}{#1}}
\newcommand{\gdsem}[1][\flp]{\mathit{DS}(#1)}
\newcommand{\gfsem}[1][\flp]{\mathit{FS}(#1)}

\newcommand{\edptext}{ed-program\xspace}
\newcommand{\edpstext}{ed-programs\xspace}

\newcommand{\Saadhead}{\ann{\literal_1}{c_1}\flpor\dots\flpor\ann{\literal_k}{c_k}}
\newcommand{\Saadbody}{\ann{\literal_{k+1}}{c_{k+1}},\dots,\ann{\literal_m}{c_m},\flpnot\ann{\literal_{m+1}}{c_{m+1}},\dots,\flpnot\ann{\literal_n}{c_n}}
\newcommand{\Saaddr}{\Saadhead\lpif{}\Saadbody}
\newcommand{\Saadleq}{\leq_A}

\newcommand{\fasptext}{FASP\xspace}
\newcommand{\faspstext}{FASPs\xspace}

\newcommand{\lkand}{\flpand_l}
\newcommand{\lkor}{\flpor_l}
\newcommand{\lkmax}{\flpor_g}
\newcommand{\lkmin}{\flpand_g}
\newcommand{\lkc}[1]{\overline{#1}}
\newcommand{\lkset}[1][k]{\lset_{#1}}
\newcommand{\lkred}[1][\fint]{\flp^{#1}}
\newcommand{\lkshift}[1][\flp]{#1^{\mathrm{sh}}}

